\begin{enumerate}
    \item Encontre o valor da expressão
    $$\sec(210^{\circ})-\tan(315^{\circ}) cossec(120^{\circ})$$
    \item Sabendo que $\cos(\theta)=- \frac{\sqrt{5}}{5}$ e que $\theta$ está no terceiro quadrante, determine os valores das outras cinco funções trigonométricas em $\theta$.
    \item Sabendo que $\tan(x)= \frac{\sqrt{2}}{4}$ e que $0 <x < 90^{\circ}$, determine $cossec(x)$ usando identidades.
    \item Calcule o valor de 
    \begin{enumerate}
        \item $\sec(-\frac{13\pi}{4})$
        \item $\sin(-\frac{\pi}{6})$
        \item $\cot(\frac{17\pi}{6})$
    \end{enumerate}
    \item Faça o esboço do gráfico das seguintes funções, destacando também o período e imagem e o deslocamento horizontal de cada uma delas.
    \begin{enumerate}
        \item $f(x)=3 \sin(x+ \frac{\pi}{2})$
        \item $g(x)=1+\frac{3}{2}\cos(\frac{x}{2})$
        \item $h(x)=\frac{1}{2}-\cos(4\pi x)$
        \item $f(x)=2-\sin(\frac{x}{2}-\frac{\pi}{2})$
    \end{enumerate}
    \item Simplifique as expressões, sempre lembrando de assumir que os denominadores são não nulos.
    \begin{enumerate}
        \item $\frac{\sin(x)}{\cos(x)+1}+\frac{\cos(x)+1}{\sin(x)}$
        \item $2\sin(x) \cot(x)-\cos(x)$
        \item $\frac{cossec(x)}{\cos(x)}+\frac{\sec(x)}{\sin(x)}$
        \item $\tan^2(x)-\tan^2(x)\sin^2(x)$
    \end{enumerate}
    \item Resolva as seguintes equações trigonométricas
    \begin{enumerate}
        \item $\sqrt{3}\tan(x)=1$
        \item $8 \cos(\frac{x}{3})+4=0$
        \item $2 \cos^2(x)-\cos(x)=0$
        \item $\tan^2(x)= \sqrt{3} \tan(x)$
        \item $\tan^2(x)=1$
        \item $7\sin(x)+4=2\sin^2(x)$
    \end{enumerate}
    \item A altura da cabine de uma roda gigante é descrita em função do tempo (em min) por
$$h(t)=76+75 \sin \left( \frac{\pi t}{15}- \frac{\pi}{2}\right)$$
\begin{enumerate}
    \item Determine as alturas máxima e mínima da cabine.
    \item Determine quanto tempo a roda gigante demora para dar uma volta completa (ou seja, o período de $h(t)$).
    \item Usando um programa gráfico, trace a curva $h(t)$.
\end{enumerate}
\end{enumerate}